\section{Goals \& Experiments}
\subsection{Goals}
The design of the rocket is based on the goals set during its definition. The goals are in order :
\begin{enumerate}
	\item Launch the rocket --- we have not launched a rocket in the past two years, so this is our main priority.
	\item Get interesting results from our two experiments --- see more details in the experiment section.
	\item Cement the club knowledge and create a solid team for the following years.
	\item Optimize the rocket --- using the knowledge of new members, we aim to optimize the rocket both aerodynamically and structurally, using CFD, FEA, and parametric and topology optimization.
	\item Optimize and predict the rocket's trajectory --- using an in-house simulator, we want to predict the trajectory as accurately as possible (verified against StabTraj) and stabilize the pitching and rolling dynamics using theoretical models, not only experimental ones (e.g.\ diameter-based static margin).
	\item Get live feedback from the experiments and the rocket state during the flight.
\end{enumerate}

\subsection{Experiments}
The following paragraphs describe the main experiments planned for the Iphigénie 60 rocket.

\paragraph{Air brake} is the main experiment. The current design status and controller approach are detailed in section \ref{sec:sys-airbrake}.
\begin{itemize}
    \item \textit{Goal:} characterize and validate the air brake's ability to control the rocket's drag so that it reaches a target apogee $h_{\textrm{target}}$, given that the motor's thrust is sized to overshoot $h_{\textrm{target}}$ by default.
    \item \textit{Parameters to study:} additional drag force (or drag coefficient) produced by the air brake as a function of deployment angle, the resulting apogee, and the vertical velocity/acceleration profile during the braking phase.
    \item \textit{Measurement range and precision:} \todo{}
    \item \textit{Hypothesis:} without the air brake, the motor drives the rocket to overshoot $h_{\textrm{target}}$; by extending drag flaps during ascent, the air brake can bleed off the excess energy and bring the achieved apogee down to $h_{\textrm{target}}$.
    \item \textit{Expected outcome:} the achieved apogee lands within a few percent of $h_{\textrm{target}}$, markedly closer than the uncontrolled (brakeless) apogee predicted in table \ref{tab:1}.
\end{itemize}

\paragraph{Reaction wheel} experiment. Background on our previous attempts is given in section \ref{sec:sys-reaction-wheel}.
\begin{itemize}
    \item \textit{Goal:} analyze whether a motor-driven reaction wheel can control the rocket's roll rate in flight, counteracting roll induced by fin misalignment and roll--pitch coupling.
    \item \textit{Parameters to study:} roll rate over time, reaction-wheel motor torque/current, and the coupling between roll rate and pitch rate.
    \item \textit{Measurement range and precision:} \todo{}
    \item \textit{Hypothesis:} a motor-driven reaction wheel can actively damp the rocket's roll rate independently of fin-induced torque, keeping it near zero during the camera-recording phase of flight despite roll--pitch coupling.
    \item \textit{Expected outcome:} the roll rate measured with the reaction wheel active is substantially lower than in our previous (uncontrolled) attempts, giving steadier camera footage.
\end{itemize}

\subsection{Additional elements}

\paragraph{Telemetry}

The onboard sensors will transmit live data to a ground station, the goal is to be able to follow the exact state of the rocket live during the flight and record the data in addition to the onboard storage.\\
We will also include cameras to record the flight and we might try to have a live video feed from the rocket.

\paragraph{Radar tracking station}

We are collaborating with a previous member of the club that is currently building a radar tracking station.\\
A turret carrying the antenna is pointed at the rocket tracking everything that is possible about it, this will then allow us to compare measurements with the onboard sensors.\\
This project is developped independently from the FUSEX and does not require any modification except a small reflector that will be placed in the nosecone thus if it is not finished in time or cannot be used during C'Space it will not affect the FUSEX.
We have more details available if needed.

\subsection{Timeline}
The project is planned according to the following milestones:
\begin{itemize}
    \item Design finished --- end of December.
    \item Optimizations finished --- mid-January.
    \item Manufacturing finished --- end of February.
\end{itemize}
